\documentclass[10pt,letterpaper]{article}
\usepackage[margin=0.72in,headheight=15pt]{geometry}
\usepackage[T1]{fontenc}
\usepackage{lmodern,microtype,amsmath,booktabs,tabularx,array,graphicx,caption,xcolor,fancyhdr,titlesec,xurl}
\usepackage[colorlinks=true,linkcolor=teal,citecolor=teal,urlcolor=teal]{hyperref}
\definecolor{ink}{HTML}{17374A}
\definecolor{teal}{HTML}{146C80}
\setlength{\parindent}{0pt}
\setlength{\parskip}{4pt}
\renewcommand{\arraystretch}{1.08}
\titleformat{\section}{\Large\bfseries\color{ink}}{\thesection}{0.6em}{}
\titleformat{\subsection}{\normalsize\bfseries\color{ink}}{}{0em}{}
\titlespacing*{\section}{0pt}{7pt}{8pt}
\titlespacing*{\subsection}{0pt}{8pt}{3pt}
\captionsetup{font=small,labelfont=bf,skip=4pt}
\pagestyle{fancy}
\fancyhf{}
\lhead{\small SPY volatility forecasting}
\rhead{\small Cumulative research report}
\cfoot{\small\thepage}
\renewcommand{\headrulewidth}{0.3pt}
\newcolumntype{L}[1]{>{\raggedright\arraybackslash}p{#1}}
\newcommand{\code}[1]{\texttt{\small #1}}
\newcommand{\sd}{\operatorname{sd}}
\input{generated/macros.tex}
\hypersetup{pdftitle={SPY Volatility Forecasting: Cumulative Research Report},pdfauthor={Pradeep Pandey}}
\begin{document}
\thispagestyle{empty}
{\small\color{teal}\textbf{INDEPENDENT STUDY / WEEKLY RESEARCH EDITION}}

{\LARGE\bfseries\color{ink} Machine Learning for Short-Term\\[3pt] SPY Volatility Forecasting}

Pradeep Pandey \textbar\ University of Wisconsin--Stout\\
Faculty supervisor: Dr. Augustine Twumasi \textbar\ \ReportDate

\section{Introduction and motivation}
Can a small regression model use information available after today's close to
forecast the variability of SPY's daily returns over the next five trading
sessions more accurately than a recent historical-volatility benchmark?

The project studies volatility regression. Its contribution is a transparent
comparison with a fixed target, a reproducible dataset, a simple benchmark, and
an explicit account of what the evaluation can establish. Price levels and
up/down classification are outside the current research question.

SPY seeks to track the S\&P 500's price and yield performance before expenses
\cite{spy}. A single broad equity instrument keeps the independent study
manageable while retaining the market variability that motivates forecasting.
The relevant information set ends at the close of session $t$; the daily return
ending at that close is available to the forecaster.

\subsection{Current research stage}
This edition integrates the requested pooled 2015--2025 comparison into the
existing Python project and verifies it offline using the original cached
Yahoo Finance snapshot. The chronological pilot remains historical context.
The newer research archive described in the handoff is unavailable, so its
results are reference values rather than an exactly reproduced run. The
\TestCount{} pooled test rows are reserved and unscored. This is a completed
development-stage comparison, not a final evaluation of future performance.

\section{Literature review}
\textbf{Measurement.} Standard deviation quantifies spread around a mean in
the units of the observations. Squared deviations make it sensitive to extremes
\cite{nist}. Normality is not required to calculate it. Here it is an interpretable,
but noisy, proxy computed from only five future daily returns. It does not
directly reveal latent conditional volatility.

\textbf{Multiple horizons.} Corsi's HAR-RV framework models volatility components
at different horizons \cite{corsi}. This motivates retaining short and longer
historical windows. Our 5-, 10-, and 20-session sample standard deviations are
computed from daily returns; this study is not a canonical HAR-RV replication.

\textbf{Model comparisons.} Breiman's Random Forests combine randomized tree
predictors \cite{breiman}, offering a nonlinear alternative to a regularized
linear model. Christensen, Siggaard, and Veliyev compare machine learning with
HAR models in a realized-variance setting \cite{christensen}. Their evidence
motivates a controlled comparison; their intraday measurement and application
differ from this five-daily-return target, so their reported gains cannot be
transferred to SPY here.

\textbf{Evaluation.} Time-series observations can remain related across randomly
selected subsets. Scikit-learn's evaluation guidance recommends considering
temporal ordering when estimating future generalization \cite{cv}. We retain the
professor's requested random split and explicitly measure its overlap.

\newpage
\section{Data and variables}
\subsection{One preserved snapshot}
The current source is Yahoo Finance through \code{yfinance}, with daily SPY,
start 1 November 2014 inclusive, end 1 January 2026 exclusive, and
\code{auto\_adjust=False}. Returns use the explicit \code{Adj Close} field.
The snapshot contains \RawCount{} observations, 3 November 2014--31 December
2025. Its original acquisition time is retained in UTC; no new prices were
downloaded for this edition. The handoff's unavailable snapshot came through
Yahoo's chart endpoint directly, not an execution of \code{yfinance}.

The calendar audit passes for all expected NYSE sessions. Dates are unique,
ordered, timezone-free session labels; prices are positive and finite. Raw
unadjusted OHLC bounds and corporate-action fields pass internal checks
(\DividendCount{} dividend events; \SplitCount{} split events). Adjusted close
is not compared with unadjusted high/low. Holidays are not missing sessions and
early closes remain observations. No absent-session prices are forward-filled;
no genuine extreme returns are removed. These checks do not independently
establish Yahoo price accuracy; cross-vendor validation has not been performed.

\subsection{Definitions and information timing}
Let $P_t$ denote adjusted close at actual trading session $t$, and
$r_t=\log(P_t/P_{t-1})$. All standard deviations use $\mathrm{ddof}=1$.
Stored values are decimals: $0.01=1\%$.

\begin{tabularx}{\linewidth}{@{}L{1.7in}L{2.5in}X@{}}
\toprule
Column & Exact definition & Role\\
\midrule
\code{log\_return} & $r_t$ & Predictor\\
\code{volatility\_5} & $\sd(r_{t-4},\ldots,r_t)$ & Predictor\\
\code{volatility\_10} & $\sd(r_{t-9},\ldots,r_t)$ & Predictor\\
\code{volatility\_20} & $\sd(r_{t-19},\ldots,r_t)$ & Predictor and benchmark\\
\code{momentum\_5} & $\sum_{j=0}^{4}r_{t-j}=\log(P_t/P_{t-5})$ & Predictor\\
\code{future\_volatility\_5d} & $\sd(r_{t+1},\ldots,r_{t+5})$ & Response only\\
\bottomrule
\end{tabularx}

In particular,
\[
 \bar r_{t,+}=\frac{1}{5}\sum_{j=1}^{5}r_{t+j},\qquad
 y_t=\sqrt{\frac{1}{4}\sum_{j=1}^{5}(r_{t+j}-\bar r_{t,+})^2}.
\]
The response is \textbf{daily-return volatility measured over five future
sessions}. It is neither annualized volatility nor the volatility of a
cumulative five-day return. The square root restores return units; the short
window and sensitivity to extremes limit its precision as a volatility proxy.

\subsection{Construction and audit metadata}
Returns and every rolling window are calculated on the full ordered price
history before random assignment. Twenty returns require 21 prices; the 2014
history provides warm-up for 2015 forecasts. Requiring complete inputs and
outcomes leaves \EligibleCount{} eligible forecasts, 2 January 2015--23 December
2025. The final five 2025 forecast dates lack complete outcomes and are excluded.

The model uses an explicit five-column predictor allowlist. The legacy response
column \code{target} retains the same meaning under its new name,
\code{future\_volatility\_5d}. The duplicate \code{baseline} is excluded from $X$.
Forecast date, original session
number, row ID, and target start/end dates remain linked metadata. Independent
NumPy slices verify all saved training and validation windows against raw prices.

\newpage
\section{Exploratory data analysis}
All figures and correlations in this section use only the \TrainingCount{}
training rows from the designated pooled run. The target median is
\TargetMedian\%, its mean \TargetMean\%, its 95th percentile \TargetUpper\%,
and its maximum \TargetMaximum\%. The right tail matters for forecast errors;
rare large values are retained as part of the research problem.

\begin{center}
\includegraphics[width=0.80\linewidth]{figures/training_distribution.pdf}
\captionof{figure}{Training target distribution in daily-volatility percent. The
full observed range is retained; the cumulative plot shows the tail's size.}
\end{center}
\begin{center}
\input{generated/correlations.tex}
\end{center}
Recent volatility is positively associated with the future target. The 5- and
10-session predictors correlate at \FeatureCorrelation{}, showing substantial
shared information. Correlation describes association, not forecast accuracy or
causation. These descriptive results motivate comparing models with the same
predictors; they do not justify an automatic feature deletion.

\begin{center}
\includegraphics[width=0.76\linewidth]{figures/training_scatter.pdf}
\captionof{figure}{Five predictors versus the response. All axes use percent
units; points are training observations. No iid significance tests are claimed.}
\end{center}

\newpage
\subsection{History and serial dependence}
The historical view places sampled observations on their actual dates. Large
return movements and elevated volatility cluster in time. Short future-window
dispersion can move sharply even when the 20-session benchmark changes more
slowly. A pooled sample contains several market conditions, but this fact alone
does not establish performance on an unseen later period.

\begin{center}
\includegraphics[width=0.88\linewidth]{figures/training_history.pdf}
\captionof{figure}{Training observations only, on actual forecast dates. Points
avoid implying that adjacent retained rows are consecutive trading sessions.}
\end{center}

\begin{center}
\includegraphics[width=0.82\linewidth]{figures/training_dependence.pdf}
\captionof{figure}{Descriptive lag correlations. A pair is included only if both
endpoints are training rows separated by the stated number of original sessions.}
\end{center}
Pair counts range from \LagPairMin{} to \LagPairMax{} across the reported lags;
the exact count is saved for every variable and lag. The randomized training
subset is never compressed into a fictitious daily sequence. Dependence and
overlapping target windows preclude ordinary independent-observation significance
claims. The target itself shares returns at short lags by construction.

\newpage
\section{Methodology}
\subsection{Requested pooled design}
The eligible rows are ordered chronologically, then permuted once using
NumPy's default random generator, seed 42:
\begin{center}\code{np.random.default\_rng(42).permutation(N)}\end{center}
The first
$\lfloor0.70N\rfloor$ rows train the models, the next
$\lfloor0.15N\rfloor$ validate them, and the remainder form the reserved test
subset. There is no stratification or seed search. This gives
\TrainingCount{} / \ValidationCount{} / \TestCount{} rows. Exact membership is
saved; sorting within subsets does not change identities.

The October 14 pilot used chronological partitions. The requested pooled
experiment is now a separate configuration built from the complete snapshot.
The handoff's single-permutation rule supersedes the older report's proposed
two-call \code{train\_test\_split} recipe, which would assign different rows.

\subsection{Overlap and prior-use limitations}
Adjacent targets share four of their five returns. The table counts held-out
rows whose future-return windows intersect any training target window or any
training feature window. It uses session metadata only, not test response values.

\begin{center}\input{generated/overlap.tex}\end{center}

Disjoint row IDs therefore do not give a leakage-free future forecasting
evaluation. Later training features also contain returns from earlier held-out
outcomes. The reported scores are a \textbf{pooled-period random-holdout
comparison}, with information overlap, rather than a demonstration of future-only
performance. Some 2015--2023 history was inspected in earlier work: creating a
new reserved subset cannot make it a never-before-seen research sample.

\subsection{Fixed algorithms and preprocessing}
The benchmark predicts $\hat y_t=\mathrm{volatility\_20}_t$. The fitted models
use the same five predictors and the same training and validation identities.

\begin{tabularx}{\linewidth}{@{}L{1.25in}X@{}}
\toprule Model & Fixed specification\\\midrule
Ridge & StandardScaler fitted only on training rows, then Ridge with $\alpha=1$.\\
Random Forest & 200 trees; maximum depth 5; minimum leaf 10; \code{max\_features="sqrt"}; seed 42; one job.\\
Gradient boosting & 200 trees; depth 2; learning rate 0.03; minimum leaf 10; squared-error loss; seed 42.\\
\bottomrule
\end{tabularx}

No tuning search is performed. Predictions have a prespecified zero floor;
\FloorCount{} validation predictions required flooring. This run uses
scikit-learn \SklearnVersion{}; full package versions and source hashes are saved.

\subsection{Evaluation}
For errors $e_i=100(\hat y_i-y_i)$, MAE is $n^{-1}\sum|e_i|$, RMSE is
$\sqrt{n^{-1}\sum e_i^2}$, and bias is $n^{-1}\sum e_i$. Units are
\textbf{daily-volatility percentage points (pp)}. Every model is scored on the
same \ValidationCount{} validation rows. High-volatility validation rows use the
training target's 95th percentile as a fixed descriptive threshold. The normal
analysis and report commands never predict or score the reserved test subset.

\newpage
\section{Results and discussion}
\subsection{Current pooled validation run}
\begin{center}\input{generated/validation.tex}\end{center}
All three fitted models improve aggregate MAE and RMSE over the benchmark in
this comparison. Gradient boosting has the lowest aggregate errors, but this
single dependent validation sample does not establish a final winner. No test
performance or statistical significance is claimed.

\begin{center}\input{generated/high_volatility.tex}\end{center}
For validation targets at or above the training 95th percentile
(\TargetUpper\% daily volatility), all learned models underpredict on average.
Lower aggregate error therefore coexists with a weakness on larger-volatility
observations. Bias and conditional errors should accompany aggregate scores.

\subsection{Reconciliation with the separate handoff analysis}
The following differences are current scores minus the handoff's six-decimal
reference scores; they are descriptive comparisons, not test tolerances.
\begin{center}\input{generated/discrepancies.tex}\end{center}
The count and split protocol agree. Baseline and Ridge are close to the quoted
values, with small tree-model differences. The cached CSV hash and acquisition
method differ from the missing archive. Different CSV schemas alone change a
hash without proving price changes; adjusted-price revisions, floating-point
inputs, and library versions may also affect fitted trees. Without the archive,
the exact contributions cannot be isolated. No data or parameters were changed
to force agreement, and exact archive reproduction is not claimed.

\subsection{Historical chronological pilot, retained separately}
\begin{center}\small\input{generated/historical.tex}\end{center}
These preserved October 14 results used 1,758 training and 496 validation rows
on the earlier chronological split. The learned pilot models were OLS and a
forest with \code{max\_features=1.0}, not the current Ridge and forest
specifications. Their archived scores are retained, not presented as a new
execution. The original repository baseline was rerun offline for this edition
and reproduces its 0.3462 MAE and 0.4664 RMSE (pp). Its fixed every-fifth check
has 100 rows, MAE 0.3577 and RMSE 0.4663. Different validation samples prevent
direct ranking of the chronological and pooled scores.

\newpage
\section{Conclusion and future work}
The project now has one integrated, offline-reproducible pooled development
experiment with explicit definitions, a verified snapshot, training-only
preprocessing and EDA, saved split identities, and a common model comparison.
The learned models reduce aggregate validation errors, while high-volatility
underprediction and cross-partition information overlap limit the conclusion.
The test subset remains reserved and unscored.

\subsection{Next research stage: temporal robustness, not yet executed}
A separately named chronological or walk-forward experiment is the next
methodological decision. Before execution, specify exact fit cutoffs, validation
dates, permitted row roles, refit frequency, and boundary exclusions. At each
fit cutoff $c$, every training label must satisfy $\mathrm{target\_end}\le c$;
evaluation forecast dates must follow the declared fitting time. If outcomes
must also be disjoint across folds, purge rows whose outcome intervals intersect
the next evaluation interval. A row-count gap alone should not replace these
timestamp checks.

The pooled reserved rows must not silently become tuning data for that study.
Because parts of the history have already informed development and pooled
training includes later years, a temporal reanalysis of this history would be
robustness evidence, not a newly untouched test. Agree its dates and roles with
the supervisor before running it. The final test remains outside routine weekly
updates; model selection, any small ablation, and the reporting protocol must be
fixed before a separately authorized final evaluation.

\subsection{Cumulative reporting and reproducibility}
One stable LaTeX source is updated as evidence accumulates. Numerical tables
are generated from saved CSV/JSON outputs, and all EDA figures refer to the
same identified pooled run. Dated and current PDFs are built from the same
source bundle. The decision log outside the main narrative records changes
without turning each weekly edition into a new research direction.

{\footnotesize
\textbf{Edition provenance.} \ReportDate{} (America/Chicago).
Pooled run: \texttt{\RunID}. Source base commit: \texttt{\SourceCommit}
with local integration changes; exact executing source hashes and working-tree
status are retained in the run record. Acquisition UTC: \texttt{\SnapshotTime}.
Snapshot SHA-256:\\
{\scriptsize\texttt{\SnapshotHash}}.
The report build record links generated tables and figures to their input hashes.
Original supplied materials are preserved separately.}

\begingroup\small
\bibliographystyle{plain}
\bibliography{references}
\endgroup
\end{document}
